% GENERATED FILE. Edit canonical lessons, then run npm run generate:books.
% canonical-source-hash: fnv1a64:9a4af4b8982cc850
% canonical-lessons: MR-C443-odha, MR-C443-kalva, MR-C443-dharan, MR-C443-sarovar, MR-C443-khalkhalne

\chapter{Streams and Dams}
\label{ch:mr-streams-and-dams}
\hlchaptermodality{443}

\begin{chapteropening}
\textbf{By the end of this chapter:} \emph{I can talk about streams, canals, dams and lakes, and the sound of running water.}

Complete the last lesson of chapter 443: I can talk about streams, canals, dams and lakes, and the sound of running water.
\end{chapteropening}

\section[\texorpdfstring{\mr{ओढा} (oḍhā)}{ओढा (oḍhā)}]{\mr{ओढा} (oḍhā) — a stream; a brook}

\label{lesson:MR-C443-odha}



\begin{quote}
Before the new one: say the Marathi for a foothill, then the Marathi for a mountain range.
\end{quote}

\subsection*{You'll want to know: \mr{ओढा}}
\textbf{\mr{ओढा}} — \emph{oḍhā} — \textquotedblleft{}a stream; a brook\textquotedblright{}.

\begin{quote}
\mr{गावाबाहेरून} \mr{एक} \mr{लहान} \mr{ओढा} \mr{वाहतो}.
\end{quote}

\begin{quote}
\emph{gāvābāherūn ek lahān oḍhā vāhto.}
\end{quote}

\begin{quote}
A small stream flows past the village.
\end{quote}

Masculine. Plural: \mr{ओढे}, \emph{oḍhe}.

Smaller than a \mr{नदी}; in the dry season many \mr{ओढे} dry up.

\begin{cousinweb}
\textbf{\mr{ओढा}} (\emph{oḍhā}) is a Marathi word for a small stream; its origin is not settled.
\end{cousinweb}

\begin{grammarlens}[title={What you've built}]
One more word for getting around.
\end{grammarlens}

\subsection*{Your turn}
\begin{itemize}
\raggedright
  \item \emph{Say it:} \emph{oḍhā}
  \item \emph{Say it:} \emph{oḍhā}, once more
  \item \emph{Say it:} \emph{gāvābāherūn ek lahān oḍhā vāhto.} — A small stream flows past the village.
  \item \emph{Recall:} say the Marathi for a foothill, then the Marathi for a mountain range, then say \emph{oḍhā} again
\end{itemize}

\begin{tcolorbox}[breakable,colback=teal!4,colframe=teal!35!black,title={Before you move on}]
What does \mr{ओढा} mean? (\textquotedblleft{}A stream; a brook\textquotedblright{}.) Say it once more.
\end{tcolorbox}
\section[\texorpdfstring{\mr{कालवा} (kālvā)}{कालवा (kālvā)}]{\mr{कालवा} (kālvā) — a canal}

\label{lesson:MR-C443-kalva}



\begin{quote}
Before the new one: say the Marathi for a mountain range, then the Marathi for a stream.
\end{quote}

\subsection*{You'll want to know: \mr{कालवा}}
\textbf{\mr{कालवा}} — \emph{kālvā} — \textquotedblleft{}a canal\textquotedblright{}.

\begin{quote}
\mr{नदीचं} \mr{पाणी} \mr{कालव्याने} \mr{शेतात} \mr{जातं}.
\end{quote}

\begin{quote}
\emph{nadīcaṁ pāṇī kālavyāne śetāt jātaṁ.}
\end{quote}

\begin{quote}
Water from the river goes to the fields by canal.
\end{quote}

Masculine.

Before an ending \mr{कालवा} becomes \mr{कालव्या}-: \mr{कालव्याने}, by the canal.

\begin{cousinweb}
\textbf{\mr{कालवा}} (\emph{kālvā}) is a Marathi word for a dug channel of water; its origin is not settled.
\end{cousinweb}

\begin{grammarlens}[title={What you've built}]
One more word for getting around.
\end{grammarlens}

\subsection*{Your turn}
\begin{itemize}
\raggedright
  \item \emph{Say it:} \emph{kālvā}
  \item \emph{Say it:} \emph{kālvā}, once more
  \item \emph{Say it:} \emph{nadīcaṁ pāṇī kālavyāne śetāt jātaṁ.} — Water from the river goes to the fields by canal.
  \item \emph{Recall:} say the Marathi for a mountain range, then the Marathi for a stream, then say \emph{kālvā} again
\end{itemize}

\begin{tcolorbox}[breakable,colback=teal!4,colframe=teal!35!black,title={Before you move on}]
What does \mr{कालवा} mean? (\textquotedblleft{}A canal\textquotedblright{}.) Say it once more.
\end{tcolorbox}
\section[\texorpdfstring{\mr{धरण} (dharaṇ)}{धरण (dharaṇ)}]{\mr{धरण} (dharaṇ) — a dam}

\label{lesson:MR-C443-dharan}



\begin{quote}
Before the new one: say the Marathi for a stream, then the Marathi for a canal.
\end{quote}

\subsection*{You'll want to know: \mr{धरण}}
\textbf{\mr{धरण}} — \emph{dharaṇ} — \textquotedblleft{}a dam\textquotedblright{}.

\begin{quote}
\mr{पावसाळ्यात} \mr{धरण} \mr{भरतं}.
\end{quote}

\begin{quote}
\emph{pāvsāḷyāt dharaṇ bhartaṁ.}
\end{quote}

\begin{quote}
In the rainy season the dam fills up.
\end{quote}

Neuter.

A noun of its own: \mr{धरण} \mr{भरणे}, a dam fills up.

\begin{cousinweb}
\textbf{\mr{धरण}} (\emph{dharaṇ}) is inherited: CDIAL lists Marathi \emph{dharaṇ}, \textquotedblleft{}dam\textquotedblright{}, under Sanskrit \emph{dharaṇa}, \textquotedblleft{}supporting, a prop\textquotedblright{} (CDIAL 6742).
\end{cousinweb}

\begin{grammarlens}[title={What you've built}]
One more word for getting around.
\end{grammarlens}

\subsection*{Your turn}
\begin{itemize}
\raggedright
  \item \emph{Say it:} \emph{dharaṇ}
  \item \emph{Say it:} \emph{dharaṇ}, once more
  \item \emph{Say it:} \emph{pāvsāḷyāt dharaṇ bhartaṁ.} — In the rainy season the dam fills up.
  \item \emph{Recall:} say the Marathi for a stream, then the Marathi for a canal, then say \emph{dharaṇ} again
\end{itemize}

\begin{tcolorbox}[breakable,colback=teal!4,colframe=teal!35!black,title={Before you move on}]
What does \mr{धरण} mean? (\textquotedblleft{}A dam\textquotedblright{}.) Say it once more.
\end{tcolorbox}
\section[\texorpdfstring{\mr{सरोवर} (sarovar)}{सरोवर (sarovar)}]{\mr{सरोवर} (sarovar) — a large lake}

\label{lesson:MR-C443-sarovar}



\begin{quote}
Before the new one: say the Marathi for a canal, then the Marathi for a dam.
\end{quote}

\subsection*{You'll want to know: \mr{सरोवर}}
\textbf{\mr{सरोवर}} — \emph{sarovar} — \textquotedblleft{}a large lake\textquotedblright{}.

\begin{quote}
\mr{सरोवराच्या} \mr{काठावर} \mr{खूप} \mr{पक्षी} \mr{असतात}.
\end{quote}

\begin{quote}
\emph{sarovarācyā kāṭhāvar khūp pakṣī astāt.}
\end{quote}

\begin{quote}
There are lots of birds on the shore of the lake.
\end{quote}

Neuter.

A literary word, bigger than a \mr{तलाव}.

\begin{cousinweb}
\textbf{\mr{सरोवर}} (\emph{sarovar}) is a tatsama, taken from Sanskrit \emph{sarovara}, \textquotedblleft{}a lake, a big pond\textquotedblright{} (Monier-Williams).
\end{cousinweb}

\begin{grammarlens}[title={What you've built}]
One more word for getting around.
\end{grammarlens}

\subsection*{Your turn}
\begin{itemize}
\raggedright
  \item \emph{Say it:} \emph{sarovar}
  \item \emph{Say it:} \emph{sarovar}, once more
  \item \emph{Say it:} \emph{sarovarācyā kāṭhāvar khūp pakṣī astāt.} — There are lots of birds on the shore of the lake.
  \item \emph{Recall:} say the Marathi for a canal, then the Marathi for a dam, then say \emph{sarovar} again
\end{itemize}

\begin{tcolorbox}[breakable,colback=teal!4,colframe=teal!35!black,title={Before you move on}]
What does \mr{सरोवर} mean? (\textquotedblleft{}A large lake\textquotedblright{}.) Say it once more.
\end{tcolorbox}
\section[\texorpdfstring{\mr{खळखळणे} (khaḷkhaḷṇe)}{खळखळणे (khaḷkhaḷṇe)}]{\mr{खळखळणे} (khaḷkhaḷṇe) — to gurgle}

\label{lesson:MR-C443-khalkhalne}



\begin{quote}
Before the new one: say the Marathi for a dam, then the Marathi for a large lake.
\end{quote}

\subsection*{You'll want to know: \mr{खळखळणे}}
\textbf{\mr{खळखळणे}} — \emph{khaḷkhaḷṇe} — \textquotedblleft{}to gurgle\textquotedblright{}.

\begin{quote}
\mr{दगडांमधून} \mr{पाणी} \mr{खळखळत} \mr{वाहतं}.
\end{quote}

\begin{quote}
\emph{dagaḍāmmadhūn pāṇī khaḷkhaḷat vāhtaṁ.}
\end{quote}

\begin{quote}
The water gurgles as it flows between the stones.
\end{quote}

Past (he …): \mr{तो} \mr{खळखळला}, \emph{to khaḷkhaḷlā}.

An echo verb: \mr{खळखळ} is the sound of a running stream.

\begin{cousinweb}
\textbf{\mr{खळखळणे}} (\emph{khaḷkhaḷṇe}) is an echo verb built on the sound \emph{khaḷkhaḷ}, running water over stones; its origin is not settled.
\end{cousinweb}

\begin{grammarlens}[title={What you've built}]
One more word for getting around.
\end{grammarlens}

\subsection*{Your turn}
\begin{itemize}
\raggedright
  \item \emph{Say it:} \emph{khaḷkhaḷṇe}
  \item \emph{Say it:} \emph{khaḷkhaḷṇe}, once more
  \item \emph{Say it:} \emph{dagaḍāmmadhūn pāṇī khaḷkhaḷat vāhtaṁ.} — The water gurgles as it flows between the stones.
  \item \emph{Recall:} say the Marathi for a dam, then the Marathi for a large lake, then say \emph{khaḷkhaḷṇe} again
\end{itemize}

\begin{tcolorbox}[breakable,colback=teal!4,colframe=teal!35!black,title={Before you move on}]
What does \mr{खळखळणे} mean? (\textquotedblleft{}To gurgle\textquotedblright{}.) Say it once more.
\end{tcolorbox}
